\documentclass[aspectratio=169]{beamer}
%--- Configuración del Tema y Colores ---
\usetheme{Madrid}
\usecolortheme{whale}
\setbeamertemplate{navigation symbols}{}
\setbeamertemplate{caption}[numbered]
%--- Paquetes ---
\usepackage[utf8]{inputenc}
\usepackage[spanish]{babel}
\usepackage{graphicx}
\usepackage{booktabs}
\usepackage{tikz}
\usepackage{fontawesome5}
\usepackage{listings}
\usepackage{amsmath}
\usepackage{amssymb}
\usetikzlibrary{shapes.geometric, arrows, positioning, calc, decorations.pathreplacing, trees}
%--- Configuración de Listings ---
\lstdefinestyle{codestyle}{
    basicstyle=\ttfamily\scriptsize,
    backgroundcolor=\color{gray!10},
    frame=single,
    breaklines=true,
    showstringspaces=false,
    columns=fullflexible,
    keywordstyle=\color{blue}\bfseries,
    commentstyle=\color{green!50!black},
    stringstyle=\color{red!70!black},
    numbers=left,
    numberstyle=\tiny\color{gray},
    numbersep=5pt
}
\lstdefinestyle{cppstyle}{
    style=codestyle,
    language=C++,
    morekeywords={vector, string, pair, set, map, queue, stack, priority_queue, unordered_map, unordered_set, deque, push_back, pop_back, emplace_back, begin, end, size, empty, front, back, top, insert, erase, find, count, lower_bound, upper_bound, auto}
}
\lstdefinestyle{pythonstyle}{
    style=codestyle,
    language=Python,
    morekeywords={def, return, for, in, range, if, else, elif, while, True, False, None, print, input, int, str, list, dict, set, append, pop, len, sorted, map, filter}
}
\lstdefinestyle{pseudostyle}{
    basicstyle=\ttfamily\scriptsize,
    backgroundcolor=\color{blue!5},
    frame=single,
    breaklines=true,
    showstringspaces=false,
    columns=fullflexible,
    mathescape=true
}
%--- Metadatos del Documento ---
\title[PM2]{Programación Matemática 2}
\subtitle{Clase 21: Programación Competitiva --- Introducción y Estructuras de Datos Básicas}
\author[Luis Alfredo Alvarado]{Prof. Luis Alfredo Alvarado Rodríguez}
\institute[ECFM - USAC]{
    Escuela de Ciencias Físicas y Matemáticas \\
    Universidad de San Carlos de Guatemala
}
\date{Primer Semestre, 2026}
\begin{document}
%--- 1. Diapositiva de Título ---
\begin{frame}
    \titlepage
\end{frame}
%--- 2. Tabla de Contenidos ---
\begin{frame}{Agenda de Hoy}
    \tableofcontents
\end{frame}
%===========================================================
% SECCIÓN 1: INTRODUCCIÓN
%===========================================================
\section{¿Qué es la Programación Competitiva?}
%--- 3 ---
\begin{frame}{Contexto: ¿Por Qué Algoritmos en un Curso de IA?}
    \textbf{En las unidades anteriores aprendimos a:}
    \begin{itemize}
        \item Usar LLMs para generar código y resolver problemas
        \item Diseñar prompts, conectar APIs, construir agentes
    \end{itemize}
    
    \vspace{0.3cm}
    \begin{block}{La Nueva Pregunta}
        Un LLM puede generar código\ldots\ pero ¿genera código \textbf{óptimo}? ¿Cómo sabemos si una solución es buena o si hay una 1000$\times$ más rápida?
    \end{block}
    
    \vspace{0.3cm}
    \begin{columns}[t]
        \column{0.48\textwidth}
        \textbf{Razones técnicas:}
        \begin{itemize}
            \item Evaluar eficiencia del código generado
            \item Optimizar pipelines de datos
            \item Los LLMs aún fallan en CP difícil
        \end{itemize}
        
        \column{0.48\textwidth}
        \textbf{Razones profesionales:}
        \begin{itemize}
            \item Entrevistas de FAANG usan CP
            \item Benchmark de IA (AlphaCode, o3)
            \item Pensamiento lógico riguroso
        \end{itemize}
    \end{columns}
\end{frame}
%--- 4 ---
\begin{frame}{Definición y Ecosistema}
    \begin{block}{Definición}
        La \textbf{Programación Competitiva} (CP) es una disciplina donde se resuelven problemas algorítmicos bien definidos bajo restricciones de \textbf{tiempo} y \textbf{memoria}, buscando soluciones correctas y eficientes.
    \end{block}
    
    \vspace{0.3cm}
    \textbf{Breve historia:}
    \begin{itemize}
        \item \textbf{1977:} Nace la ACM-ICPC (hoy ICPC) --- competencia universitaria más prestigiosa
        \item \textbf{1989:} Inicia la IOI (Olimpiada Internacional de Informática)
        \item \textbf{2000s:} Plataformas online: Codeforces, TopCoder, LeetCode, AtCoder
        \item \textbf{2020s:} LLMs compiten --- AlphaCode (2022), o3 rating $\sim$1800+ en Codeforces
    \end{itemize}
    
    \vspace{0.2cm}
    \begin{exampleblock}{Plataformas clave}
        \textbf{Codeforces} (contests semanales), \textbf{LeetCode} (entrevistas), \textbf{CSES} (300 problemas clásicos), \textbf{AtCoder} (comunidad JP)
    \end{exampleblock}
\end{frame}
%--- 5 ---
\begin{frame}{Anatomía de un Problema de CP}
    \centering
    \begin{tikzpicture}[scale=0.85, transform shape]
        \tikzset{
            component/.style={rectangle, rounded corners, minimum width=10cm, minimum height=0.9cm, draw, font=\scriptsize, align=left}
        }
        
        \node[component, fill=red!20] (title) at (0,3) {\textbf{Nombre del Problema:} Título descriptivo + ID};
        \node[component, fill=blue!20] (desc) at (0,1.8) {\textbf{Enunciado:} Descripción narrativa del problema a resolver};
        \node[component, fill=yellow!20] (input) at (0,0.6) {\textbf{Formato de Entrada:} Especificación exacta de cómo se reciben los datos};
        \node[component, fill=green!20] (output) at (0,-0.6) {\textbf{Formato de Salida:} Especificación exacta de lo que se debe imprimir};
        \node[component, fill=purple!20] (limits) at (0,-1.8) {\textbf{Restricciones:} $1 \leq n \leq 10^5$, tiempo: 2s, memoria: 256 MB};
        \node[component, fill=orange!20] (examples) at (0,-3) {\textbf{Casos de Ejemplo:} Pares entrada/salida para verificar comprensión};
        
        \draw[decorate, decoration={brace, amplitude=5pt}] (5.5,-1.4) -- (5.5,-3.4) node[midway, right=0.3cm, font=\tiny, text width=2.5cm] {Claves para elegir\\el algoritmo};
    \end{tikzpicture}
\end{frame}
%--- 6 ---
\begin{frame}[fragile]{Ejemplo: Un Problema Clásico}
    \begin{block}{Problema: Suma de Pares}
        Dado un arreglo de $n$ enteros y un valor objetivo $k$, determinar si existen dos elementos cuya suma sea exactamente $k$.
    \end{block}
    
    \begin{columns}
        \column{0.45\textwidth}
        \textbf{Entrada:}
\begin{lstlisting}[style=pseudostyle]
5 9
2 7 11 15 1
\end{lstlisting}
        \textit{$n=5$, $k=9$, arreglo: [2, 7, 11, 15, 1]}
        
        \vspace{0.2cm}
        \textbf{Salida:}
\begin{lstlisting}[style=pseudostyle]
YES
\end{lstlisting}
        \textit{Porque $2 + 7 = 9$}
        
        \column{0.50\textwidth}
        \textbf{¿Cómo lo resolverías?}
        \begin{itemize}
            \item \textbf{Fuerza bruta:} Todos los pares $\rightarrow O(n^2)$
            \item \textbf{Ordenar + dos punteros:} $\rightarrow O(n \log n)$
            \item \textbf{Hash set:} $\rightarrow O(n)$
        \end{itemize}
        
        \vspace{0.2cm}
        La \textbf{elección de la estructura de datos} determina la eficiencia. Veremos las tres soluciones en esta clase.
    \end{columns}
\end{frame}
%===========================================================
% SECCIÓN 2: SETUP
%===========================================================
\section{Setup: Lenguaje y Plantilla Base}
%--- 7 ---
\begin{frame}{Lenguajes Populares en CP}
    \begin{table}[]
        \centering
        \renewcommand{\arraystretch}{1.3}
        \scriptsize
        \begin{tabular}{p{2cm} p{3.5cm} p{3.5cm} p{2.5cm}}
            \toprule
            \textbf{Lenguaje} & \textbf{Ventajas} & \textbf{Desventajas} & \textbf{Uso en CP} \\
            \midrule
            \textbf{C++} & Rápido, STL poderoso & Verbose, errores crípticos & $\sim$80\% competidores \\
            \textbf{Python} & Legible, prototipado rápido & Lento ($\sim$50x más que C++) & Problemas sin TLE \\
            \textbf{Java} & Robusto, BigInteger & Verbose, más lento que C++ & $\sim$10\% competidores \\
            \bottomrule
        \end{tabular}
    \end{table}
    
    \vspace{0.2cm}
    \begin{block}{Recomendación para este curso}
        Usaremos \textbf{C++} como lenguaje principal (estándar en competencias) y \textbf{Python} para prototipado rápido y verificación.
    \end{block}
\end{frame}

%===========================================================
% SECCIÓN 3: ARREGLOS Y STRINGS
%===========================================================
\section{Arreglos, Strings y Técnicas Clave}
%--- 9 ---
\begin{frame}{Arreglos (Arrays)}
    \begin{block}{Definición}
        Secuencia contigua de elementos del mismo tipo, accesibles por índice en $O(1)$.
    \end{block}
    
    \centering
    \begin{tikzpicture}[scale=0.9]
        \foreach \i/\v in {0/3, 1/7, 2/1, 3/9, 4/4, 5/2, 6/8, 7/5} {
            \draw[fill=blue!20] (\i*1.2, 0) rectangle (\i*1.2+1, 0.8);
            \node at (\i*1.2+0.5, 0.4) {\textbf{\v}};
            \node[font=\tiny, gray] at (\i*1.2+0.5, -0.3) {\i};
        }
        \node[font=\scriptsize, left] at (-0.3, 0.4) {valor};
        \node[font=\scriptsize, left, gray] at (-0.3, -0.3) {índice};
    \end{tikzpicture}
    
    \vspace{0.3cm}
    \begin{table}[]
        \centering
        \scriptsize
        \renewcommand{\arraystretch}{1.2}
        \begin{tabular}{l c l}
            \toprule
            \textbf{Operación} & \textbf{Complejidad} & \textbf{Nota} \\
            \midrule
            Acceso por índice & $O(1)$ & Lo mejor de arreglos \\
            Búsqueda (sin orden) & $O(n)$ & Recorrido lineal \\
            Búsqueda (ordenado) & $O(\log n)$ & Búsqueda binaria \\
            Inserción al final & $O(1)$ amortizado & \texttt{push\_back} en C++ \\
            Inserción en medio & $O(n)$ & Requiere desplazar elementos \\
            \bottomrule
        \end{tabular}
    \end{table}
\end{frame}

%--- 12 ---
\begin{frame}[fragile]{Strings: Operaciones Esenciales}
    Los strings son arreglos de caracteres con operaciones especiales.
    
    \begin{columns}
        \column{0.48\textwidth}
        \textbf{Operaciones frecuentes en CP:}
        \begin{itemize}
            \item Comparación lexicográfica
            \item Búsqueda de subcadena
            \item Palíndromos
            \item Conteo de frecuencias
        \end{itemize}
        
        \vspace{0.2cm}
        \begin{alertblock}{Cuidado}
            Concatenar strings en un ciclo es $O(n^2)$. Usar \texttt{stringstream} (C++) o \texttt{join} (Python).
        \end{alertblock}
        
        \column{0.48\textwidth}
\begin{lstlisting}[style=cppstyle, numbers=none]
string s = "abracadabra";

// Frecuencia de caracteres
int freq[26] = {};
for(char c : s)
    freq[c - 'a']++;

// Verificar palindromo
bool esPalin(string &s){
    int n = s.size();
    for(int i = 0; i < n/2; i++)
        if(s[i] != s[n-1-i])
            return false;
    return true;
}
\end{lstlisting}
    \end{columns}
\end{frame}

%--- 19 ---
\begin{frame}
    \centering
    \Huge \textbf{¿Preguntas?}
    
    \vspace{0.5cm}
    \normalsize
    Prof. Luis Alfredo Alvarado Rodríguez\\
    \small Escuela de Ciencias Físicas y Matemáticas
    
    \vspace{0.5cm}
    \footnotesize
    \textbf{Ejercicios:} Two Sum (LeetCode \#1), Valid Parentheses (LeetCode \#20),\\
    Range Sum Query (LeetCode \#303), CSES Missing Number\\[0.3cm]
    \textbf{Recursos:} Laaksonen, \textit{Competitive Programmer's Handbook} (gratis online)\\
    Halim \& Halim, \textit{Competitive Programming 4}
    
    \vspace{0.3cm}
    \scriptsize
    \textit{``First, solve the problem. Then, write the code.''\\--- John Johnson}
\end{frame}
\end{document}